\section{Results}

\subsection{Category-Dependent ECE Variation (h-e1: PASS)}

ANOVA revealed highly significant variation (F=8.45, p=0.00012):

\begin{table}[h]
\centering
\begin{tabular}{@{}lc@{}}
\toprule
\textbf{Cluster} & \textbf{ECE} \\
\midrule
Finance/Economics & 0.152 \\
Science/Technology/Math & 0.169 \\
Health/Nutrition/Psychology & 0.184 \\
History/Geography/Culture & 0.193 \\
Law/Politics/Government & 0.216 \\
Religion/Philosophy/Ethics & 0.228 \\
Misconceptions/Myths/Superstitions & 0.251 \\
\bottomrule
\end{tabular}
\caption{ECE by cluster, sorted by calibration quality.}
\label{tab:ece}
\end{table}

ECE range: 0.099 (nearly double target threshold of 0.05).

\subsection{Distinct Confidence Distributions (h-m1: PASS)}

17/21 cluster pairs showed significant differences at p<0.05, exceeding threshold of 11/21. Confidence range: 0.433 (from mean 0.412 to 0.845).

\textbf{Important methodological caveat:} Due to GPU constraints, h-m1 used simulated confidence data calibrated to match h-e1 ECE patterns. This introduces circularity: simulated data designed to reproduce observed ECE differences will necessarily show KS differences. The 17/21 significant pairs result therefore demonstrates internal consistency of the simulation rather than an independent empirical finding. We retain h-m1 results for completeness but note that the h-m2 uniformity finding (Section~\ref{sec:temperature}) --- which used actual model outputs --- provides the primary mechanistic evidence.

\subsection{Temperature Uniformity (h-m2: FAIL)}
\label{sec:temperature}

All clusters converged to identical optimal temperature:

\begin{table}[h]
\centering
\begin{tabular}{@{}lc@{}}
\toprule
\textbf{Cluster} & \textbf{Optimal T} \\
\midrule
All 7 clusters & 10.0 \\
\bottomrule
\end{tabular}
\caption{Optimal temperature by cluster.}
\label{tab:temperature}
\end{table}

CV = 0.0, Range = 0.0 (targets: $>0.1$ and $>0.3$ respectively).

The uniform T=10.0 convergence indicates severe, uniform overconfidence across all categories.

\textbf{Interpretive note:} T=10.0 is the upper bound of our search range $[0.1, 10.0]$. This bound saturation creates interpretive uncertainty: either (a) the true optimal is T=10.0, or (b) the true optimal lies beyond T=10 and cluster differentiation may emerge at higher temperatures. Consequently, we can only conclude that no cluster differentiation exists \emph{within} the tested range. Future work should explore extended temperature bounds (e.g., T up to 50 or 100) to determine whether cluster-specific optima emerge at higher temperatures.
